\documentclass[11pt,a4paper]{article}

% Paquetes esenciales
\usepackage{cmap} % texto copiable/legible por ATS; debe ir antes de fontenc
\usepackage[utf8]{inputenc}
\usepackage[T1]{fontenc}
\usepackage{lmodern}
\usepackage[spanish]{babel}
\usepackage[left=1.5cm, right=1.5cm, top=1.5cm, bottom=1.5cm]{geometry}
\usepackage{hyperref}
\usepackage{enumitem}
\usepackage{titlesec}
\usepackage{xcolor}

% Configuración de enlaces
\hypersetup{
    colorlinks=true,
    linkcolor=blue,
    filecolor=magenta,
    urlcolor=blue,
    pdftitle={CV - Breyner Ismael Ciro Otero},
}

% Formato de secciones (Línea debajo del título)
\titleformat{\section}
  {\large\bfseries\uppercase}
  {}
  {0em}
  {}[\titlerule]
\titlespacing*{\section}{0pt}{1.5ex plus 1ex minus .2ex}{1ex plus .2ex}

% Desactivar numeración de páginas
\pagestyle{empty}

\begin{document}

% --- ENCABEZADO ---
\begin{center}
    {\huge \textbf{BREYNER ISMAEL CIRO OTERO}} \\ \vspace{0.15cm}
    {\large Ingeniero de Software FullStack $\cdot$ TypeScript $\cdot$ Vue/React $\cdot$ Node.js $\cdot$ GCP/AWS $\cdot$ IA aplicada} \\ \vspace{0.15cm}
    Bogotá, Colombia | +57 3183420438 | \href{mailto:breynerismael@gmail.com}{breynerismael@gmail.com} \\ \vspace{0.1cm}
    \href{https://www.linkedin.com/in/breyner-ciro-38011330b/}{LinkedIn} | \href{https://github.com/breynerciro}{GitHub} | \href{https://breynerciro.github.io/}{Portafolio}
\end{center}

\vspace{0.2cm}

% --- PERFIL PROFESIONAL ---
\section*{Perfil Profesional}
Ingeniero de Sistemas y Computación (Universidad Nacional de Colombia) que construye producto SaaS de punta a punta: frontend en Vue y React con TypeScript, backends en Node.js y Python, y datos en PostgreSQL, Firestore y BigQuery sobre GCP, Supabase y Cloudflare. Hoy construye agentes de IA (Claude, Gemini, MCP) sobre datos reales de negocio en una plataforma logística y desarrolla dos productos propios: un POS con facturación electrónica y una plataforma para consultorios odontológicos. Trabaja con pruebas automatizadas y despliegues en contenedores.

% --- HABILIDADES TÉCNICAS ---
\section*{Habilidades Técnicas}
\begin{itemize}[leftmargin=*, noitemsep]
    \item \textbf{Lenguajes:} TypeScript, JavaScript (ES6+), Python, SQL.
    \item \textbf{Frontend \& Mobile:} Vue 3 (Pinia), React, React Native, Expo, Tailwind CSS, Vite, PWA, Astro.
    \item \textbf{Backend \& APIs:} Node.js, Express, FastAPI, Django, APIs RESTful, Webhooks, SSE, Socket.IO, BullMQ, Supabase Edge Functions (Deno), Gmail API, Microsoft Graph, WhatsApp Cloud API.
    \item \textbf{IA aplicada:} Agentes sobre LLMs (Claude vía API de Anthropic y Vertex AI, Gemini), Model Context Protocol (MCP), salidas estructuradas con Zod, extracción de datos de documentos.
    \item \textbf{Datos \& Cloud:} PostgreSQL (Row-Level Security), Supabase, Firebase (Firestore, Cloud Functions), BigQuery, Cloud Run, Cloudflare Workers, MySQL, MongoDB, Redis, AWS, Docker, Nginx, Caddy, GitHub Actions.
    \item \textbf{Calidad \& Prácticas:} Vitest, Playwright, Cypress, k6, JMeter, observabilidad (OpenTelemetry, Prometheus, Sentry), Code Reviews, GitFlow, Scrum, arquitectura multitenant y de microservicios.
    \item \textbf{Idiomas:} Español (nativo), Inglés (B1), Portugués (B1).
\end{itemize}

% --- EXPERIENCIA PROFESIONAL ---
\section*{Experiencia Profesional}

\noindent
\textbf{Kila} | \textit{Desarrollador Junior de Ingeniería de Software} \hfill Octubre, 2026 – Actualidad \\
\textit{SaaS logístico para importadores y freight forwarders.}
\begin{itemize}[leftmargin=*, noitemsep, topsep=2pt]
    \item Formo parte del equipo de ingeniería de una plataforma SaaS que centraliza la operación de comercio exterior de sus clientes; aporto de punta a punta en frontend, backend serverless y datos (\textbf{Vue 3, Pinia, TypeScript, Firebase, BigQuery, Cloud Run}).
    \item Construyo agentes de IA en producto con \textbf{Claude} (API de Anthropic y Vertex AI) y \textbf{Gemini}: un asistente que arma dashboards con datos en vivo de BigQuery y un agente de soporte de primer nivel que revisa tickets y deja borradores para aprobación humana.
    \item Doy a los agentes acceso a los datos del negocio con \textbf{Model Context Protocol (MCP)}, contratos \textbf{Zod} y respuestas en streaming (SSE), aislando los datos de cada cliente con pruebas de contrato multitenant que bloquean el despliegue en CI.
    \item Fortalezco la confiabilidad y la oportunidad de los datos operativos y aduaneros con los que los clientes toman decisiones; reduje a la mitad el tiempo de sincronización con el data warehouse (de más de 2 min a cerca de 1 min).
    \item Entrego funcionalidades de operación como la conciliación de facturas contra órdenes de compra, campos personalizados para bodega y menciones en el correo integrado (\textbf{Gmail, Microsoft Graph}), con pruebas automatizadas en 4 servicios.
\end{itemize}

\vspace{0.2cm}

\noindent
\textbf{CardioDiagnostico del Sur} | \textit{Desarrollador Frontend} \hfill Enero, 2024 – Agosto, 2024
\begin{itemize}[leftmargin=*, noitemsep, topsep=2pt]
    \item Construí el sitio institucional de una IPS de cardiología como Single Page Application con \textbf{React} y \textbf{TypeScript} (servicios, estados financieros y contacto), con código fuertemente tipado.
    \item Maqueté una interfaz 100\% responsiva con \textbf{Tailwind CSS} y gestioné el paso a producción con \textbf{Vite} y despliegue en Hostinger.
    \item Facilité el agendamiento de citas con un botón de WhatsApp (mensaje prellenado) y un formulario de contacto con backend propio en \textbf{Node.js/Express} que envía los mensajes por correo.
\end{itemize}

\vspace{0.2cm}

\noindent
\textbf{Proyectos Corporativos (Freelance / U. Nacional)} | \textit{Desarrollador de Software} \hfill 2022 – 2024
\begin{itemize}[leftmargin=*, noitemsep, topsep=2pt]
    \item Integré SPAs con backends relacionales y validé su comportamiento bajo concurrencia con pruebas de carga y estrés.
    \item Estandaricé los entornos de desarrollo y producción con configuraciones reproducibles en \textbf{Docker} y docker-compose.
    \item Automaticé los flujos de integración y entrega continua con canalizaciones en \textbf{GitHub Actions}.
\end{itemize}

% --- PROYECTOS DESTACADOS ---
\section*{Proyectos Destacados}

\noindent
\textbf{BREVENTA – SaaS de Punto de Venta y Facturación Electrónica (producto propio)} \hfill 2026
\begin{itemize}[leftmargin=*, noitemsep, topsep=2pt]
    \item Desarrollo un POS multiempresa en monorepo: PWA en \textbf{Vue 3}, \textbf{Pinia} y \textbf{Tailwind CSS}, API en \textbf{Node.js}, \textbf{Express} y \textbf{TypeScript}, y un microservicio en \textbf{FastAPI} para la facturación electrónica ante la DIAN (CUFE, firma XAdES), sobre \textbf{PostgreSQL} con Row-Level Security por empresa, \textbf{Redis}, BullMQ y Socket.IO.
    \item Preparé la operación con \textbf{Docker Compose} y Caddy (HTTPS), observabilidad con OpenTelemetry, Prometheus y Sentry, 62 pruebas \textbf{Playwright} de seguridad, responsividad y accesibilidad, pruebas de carga con \textbf{k6} (50 empresas) y análisis de seguridad con semgrep, gitleaks y trivy.
\end{itemize}

\vspace{0.2cm}

\noindent
\textbf{DENTY – Gestión de Consultorios Odontológicos (producto propio, en producción)} \hfill 2026
\begin{itemize}[leftmargin=*, noitemsep, topsep=2pt]
    \item Construí una PWA en \textbf{Vue 3}, \textbf{TypeScript} y \textbf{Tailwind CSS} sobre \textbf{Supabase} (PostgreSQL con Row-Level Security, Auth, Storage, Edge Functions en Deno, pg\_cron) y \textbf{Cloudflare Workers}: pacientes, historia clínica con odontograma, agenda, presupuestos y consentimientos firmados en PDF con hash SHA-256.
    \item Integré un asistente de \textbf{WhatsApp} con \textbf{Gemini} que confirma, agenda y reagenda citas, y dejé en la base de datos las reglas del negocio: evoluciones firmadas inmutables, bitácora de auditoría y acceso solo por invitación.
\end{itemize}

\vspace{0.2cm}

\noindent
\textbf{HEMODENT – Aplicación Móvil Educativa en Salud (Proyecto de Tesis)} \hfill 2026
\begin{itemize}[leftmargin=*, noitemsep, topsep=2pt]
    \item Desarrollé una app multiplataforma (iOS/Android) para pacientes con hemofilia con \textbf{React Native}, \textbf{TypeScript} y \textbf{Expo} (SDK 57, Expo Router), con persistencia offline (AsyncStorage), animaciones (Reanimated) y recordatorios mediante notificaciones locales.
    \item Diseñé una interfaz (UX/UI) adaptativa por rango de edad, con acceso a protocolos de emergencia dental y un centro de conocimiento médico especializado.
\end{itemize}

\vspace{0.2cm}

\noindent
\textbf{OSP – Overwatch Sentinel Platform (Plataforma Edge de Visión Artificial, proyecto en equipo)} \hfill 2025
\begin{itemize}[leftmargin=*, noitemsep, topsep=2pt]
    \item Desarrollé una plataforma de vigilancia con visión por computadora usando \textbf{TensorFlow Lite} en \textbf{Raspberry Pi} (Edge Computing) para detección de objetos y optimización de grabación de video.
    \item Diseñé una arquitectura de microservicios con backend (\textbf{Python/Node.js}), frontend (\textbf{React}) y bases de datos (\textbf{PostgreSQL, MongoDB}), comunicados mediante RESTful APIs, OAuth2 (Google) y JWT.
    \item Orquesté el despliegue en 5 entornos con \textbf{Docker}, un API Gateway (\textbf{Vert.x}) con patrón Gatekeeper y proxy reverso (\textbf{Nginx}); validé alta disponibilidad y tolerancia a fallos con pruebas de estrés en \textbf{JMeter}.
\end{itemize}

% --- EDUCACIÓN ---
\section*{Educación}

\noindent
\textbf{Ingeniero de Sistemas y Computación} | Universidad Nacional de Colombia \hfill Graduado: Septiembre 2026 \\
\textit{Bogotá D.C., Colombia.}

\vspace{0.2cm}

\noindent
\textbf{Técnico Laboral por Competencias en Sistemas Informáticos} | POLINTEC \hfill Diciembre 2016 \\
\textit{Puerto Asís, Putumayo.}

% --- CERTIFICACIONES ---
\section*{Certificaciones}
\begin{itemize}[leftmargin=*, noitemsep]
    \item \textbf{Diplomado en Arquitectura y Calidad de Software (192h, 3 módulos):} Universidad Nacional de Colombia. Terminados: Ingeniería de Software Basada en la Nube y Calidad de Software. En curso: Arquitectura de Software a Gran Escala (finaliza en noviembre de 2026).
    \item \textbf{AWS Academy Graduate - Cloud Foundations (20h):} Credly (Noviembre 2025).
    \item \textbf{Curso Básico de Power BI (8h):} Datly Educación (Julio 2025).
\end{itemize}

\end{document}
